\documentclass{article}

\usepackage{amsmath,amssymb}
\usepackage{xurl}
\usepackage[hidelinks]{hyperref}
\setlength{\emergencystretch}{2em}

% Shared commands for notes in docs/paper-gaps/.

\DeclareUnicodeCharacter{2080}{\ifmmode{}_{0}\else\textsubscript{0}\fi}
\DeclareUnicodeCharacter{2081}{\ifmmode{}_{1}\else\textsubscript{1}\fi}
\DeclareUnicodeCharacter{2082}{\ifmmode{}_{2}\else\textsubscript{2}\fi}
\DeclareUnicodeCharacter{2099}{\ifmmode{}_{n}\else\textsubscript{n}\fi}

\newcommand{\Lean}{\textsc{Lean}}
\ExplSyntaxOn
\NewDocumentCommand{\leanid}{m}{
  \ifmmode
    \textsf{\detokenize{#1}}
  \else
    \group_begin:
    \urlstyle{sf}
    \tl_set:Nn \l_tmpa_tl {#1}
    \tl_replace_all:Nnn \l_tmpa_tl {\_} {_}
    \exp_args:NV \path \l_tmpa_tl
    \group_end:
  \fi
}
\ExplSyntaxOff
\providecommand{\tr}{\operatorname{tr}}
\newcommand{\tnleanIssueBase}{https://github.com/LionSR/TNLean/issues/}
\newcommand{\tnleanPrBase}{https://github.com/LionSR/TNLean/pull/}
\newcommand{\tnleanDocBase}{https://sirui-lu.com/TNLean/docs/}
\newcommand{\ghissue}[1]{\href{\tnleanIssueBase#1}{GitHub issue \##1}}
\newcommand{\ghpr}[1]{\href{\tnleanPrBase#1}{PR \##1}}
\newcommand{\doclink}[1]{\href{\tnleanDocBase#1.html}{\path{#1}}}

% Dirac notation and the identity operator, as in blueprint/src/macros/common.tex.
% \providecommand keeps note-local definitions authoritative.
\usepackage{dsfont}
\providecommand{\ket}[1]{|#1\rangle}
\providecommand{\bra}[1]{\langle #1|}
\providecommand{\braket}[2]{\langle #1 | #2 \rangle}
\providecommand{\ketbra}[2]{|#1\rangle\!\langle #2|}
\providecommand{\Id}{\mathds{1}}
\providecommand{\one}{\mathds{1}}
\providecommand{\C}{\mathbb{C}}
\providecommand{\MN}[1]{M_{#1}(\C)}
\providecommand{\MD}{M_D(\C)}

% External referencing for paper sources.  The build helper writes sanitized
% auxiliary files under build/paper-aux/ and excludes duplicated source labels.
\usepackage{xr}

% Import a generated paper auxiliary file with a caller-supplied label prefix.
% The candidate paths support builds from docs/paper-gaps/, the repository root,
% and build/paper-aux/ itself.
\newcommand{\importpaperaux}[2]{%
  \IfFileExists{../../build/paper-aux/#2.aux}{%
    \externaldocument[#1]{../../build/paper-aux/#2}%
  }{%
    \IfFileExists{build/paper-aux/#2.aux}{%
      \externaldocument[#1]{build/paper-aux/#2}%
    }{%
      \IfFileExists{#2.aux}{%
        \externaldocument[#1]{#2}%
      }{}%
    }%
  }%
}

% General external reference macros
\newcommand{\paperref}[2]{\ref{#1-#2}}
\newcommand{\papereqref}[2]{(\ref{#1-#2})}

% CPSV16 source (1606.00608 / MPDO-22-12-17-2.tex) external document and shortcuts
\importpaperaux{cpsv16-}{1606.00608/MPDO-22-12-17-2-tnlean}
\newcommand{\cpsvref}[1]{\paperref{cpsv16}{#1}}
\newcommand{\cpsveqref}[1]{\papereqref{cpsv16}{#1}}

% Verdict marker: \gapnote{<kind>}{<status>}. Kind is the policy
% classification (clarification, local-correction, scope-restriction,
% unfaithful, false-source, open-gap); status is open, wip (elimination
% underway), resolved, or historical. Machine-read by the site index;
% prints a verdict line.
\newcommand{\gapnote}[2]{\par\noindent\textbf{Verdict:} #1 (#2).\par}

\providecommand{\leanid}[1]{\path{#1}}

\title{Injective PEPS Gauge Consistency Needs a Connectivity Hypothesis}
\author{}
\date{2026-06-06}

\begin{document}
\maketitle
\gapnote{false-source}{resolved}

\section{The assertion}

The gauge-consistency step in the Fundamental Theorem for injective PEPS
asserts that two vertex-injective tensors on a simple graph which generate the
same state are related by one global family of edge gauges
\cite{Molnar2018NormalPEPS}. After the source equation labelled
\path{eq:inj_equal_edge}, blocking one vertex against its complement gives a
scalar proportionality
\begin{align}
    A_v
        &= \lambda_v\,\widetilde B_v.
    \label{eq:peps_gaugeConsistency_connectivity_gap_vertex_proportionality}
\end{align}
The source then states that ``the constants $\lambda_v$ can be incorporated
into the gauge transformations''
\cite{Molnar2018NormalPEPS}.\footnote{Repository source:
\path{Papers/1804.04964/paper_normal.tex:1207}.}

For a graph $G=(V,E)$, let $X=(X_e)_{e\in E}$ be one global family of
invertible edge gauges, and let
$\operatorname{gaugeVertex}(A,X)_v$ denote its oriented action on the tensor
$A_v$. The formal gauge-consistency statement requires
\begin{align}
    B_v(\eta)_\sigma
        &= \operatorname{gaugeVertex}(A,X)_v(\eta)_\sigma
    \label{eq:peps_gaugeConsistency_connectivity_gap_formal_gauge_consistency}
\end{align}
for every $v$, virtual configuration $\eta$, and physical index $\sigma$.
This statement is formalized by
\leanid{TNLean.PEPS.gaugeConsistency}. The headline declaration
\leanid{TNLean.PEPS.fundamentalTheorem_PEPS} reduces to this gauge-consistency
step.

\section{Failure without connectivity}

The reduction to the vertex scalars in
Eq.~\ref{eq:peps_gaugeConsistency_connectivity_gap_vertex_proportionality} is
correct. Their absorption into edge gauges is not valid on an arbitrary graph.

Fix an orientation of each edge. To absorb the nonzero vertex scalars
$\lambda_v$, one needs nonzero edge scalars $s=(s_e)_{e\in E}$. Define the
oriented endpoint scalar by
\begin{align}
    s_{v,e}
        &=
        \begin{cases}
            s_e, & \text{if $v$ is the tail of $e$},\\
            s_e^{-1}, & \text{if $v$ is the head of $e$}.
        \end{cases}
    \label{eq:peps_gaugeConsistency_connectivity_gap_oriented_edge_scalar}
\end{align}
Absorption requires the incidence equations
\begin{align}
    \prod_{e\ni v}s_{v,e}
        &= \lambda_v^{-1}
        \qquad\text{for every }v\in V.
    \label{eq:peps_gaugeConsistency_connectivity_gap_incidence_equations}
\end{align}
For every connected component $C\subseteq V$, multiplying the left-hand sides
of Eq.~\ref{eq:peps_gaugeConsistency_connectivity_gap_incidence_equations}
cancels each edge scalar against its inverse:
\begin{align}
    \prod_{v\in C}\prod_{e\ni v}s_{v,e}
        &= 1.
    \label{eq:peps_gaugeConsistency_connectivity_gap_component_cancellation}
\end{align}
Thus the reciprocal vertex-scalar family is absorbable only if
\begin{align}
    \prod_{v\in C}\lambda_v^{-1}
        &= 1
        \qquad\text{for every connected component }C.
    \label{eq:peps_gaugeConsistency_connectivity_gap_component_product_condition}
\end{align}
State equality forces only the product over all vertices. On a connected graph,
this is the single condition in
Eq.~\ref{eq:peps_gaugeConsistency_connectivity_gap_component_product_condition};
on a disconnected graph, it does not force the condition separately on each
component.

\subsection{Empty-graph witness}

Let $G=(\{0,1\},\varnothing)$, take physical dimension one, and define
bond-dimension-one tensors by
\begin{align}
    A_0
        &= 2,
    \label{eq:peps_gaugeConsistency_connectivity_gap_witness_A_zero}\\
    A_1
        &= 3,
    \label{eq:peps_gaugeConsistency_connectivity_gap_witness_A_one}\\
    B_0
        &= 6,
    \label{eq:peps_gaugeConsistency_connectivity_gap_witness_B_zero}\\
    B_1
        &= 1.
    \label{eq:peps_gaugeConsistency_connectivity_gap_witness_B_one}
\end{align}
Every vertex tensor is a single nonzero number, so both PEPS are
vertex-injective. Their states agree because
\begin{align}
    \langle A\rangle
        &= A_0A_1
         = 6,
    \label{eq:peps_gaugeConsistency_connectivity_gap_witness_state_A}\\
    \langle B\rangle
        &= B_0B_1
         = 6.
    \label{eq:peps_gaugeConsistency_connectivity_gap_witness_state_B}
\end{align}
All bond dimensions are vacuously positive because the graph has no edges.

The empty graph admits only the empty edge-gauge family. Its action is trivial:
\begin{align}
    \operatorname{gaugeVertex}(A,X)_v
        &= A_v
        \qquad\text{for every }v\in\{0,1\}.
    \label{eq:peps_gaugeConsistency_connectivity_gap_empty_gauge_action}
\end{align}
The conclusion in
Eq.~\ref{eq:peps_gaugeConsistency_connectivity_gap_formal_gauge_consistency}
would therefore imply $A=B$, but
\begin{align}
    A_0
        &= 2
         \neq 6
         = B_0.
    \label{eq:peps_gaugeConsistency_connectivity_gap_witness_contradiction}
\end{align}
Hence \leanid{TNLean.PEPS.gaugeConsistency} is false on this input when no
connectivity hypothesis is imposed.

This counterexample has been checked in \Lean. The negation of the conclusion
of \leanid{TNLean.PEPS.gaugeConsistency} is proved for this input using only
the standard logical dependencies \leanid{propext}, \leanid{Classical.choice},
and \leanid{Quot.sound}. The same input refutes
\leanid{TNLean.PEPS.fundamentalTheorem_PEPS}, because the empty edge gauge
satisfies
\begin{align}
    \operatorname{GaugeEquiv}A\,B
        &\Longleftrightarrow A=B.
    \label{eq:peps_gaugeConsistency_connectivity_gap_empty_graph_gauge_equiv}
\end{align}

\section{Corrected hypothesis}

The theorem in Ref.~\cite{Molnar2018NormalPEPS} is intended for connected
PEPS. A PEPS on a disconnected graph is a product of independent states and
lies outside the local blocking argument used there. The faithful correction
adds connectivity to \leanid{TNLean.PEPS.gaugeConsistency} and
\leanid{TNLean.PEPS.fundamentalTheorem_PEPS}, for example through
\leanid{G.Connected}, or through \leanid{G.Preconnected} together with
nonemptiness of $V$.

Under connectivity, state equality and state nonvanishing give
\begin{align}
    \prod_{v\in V}\lambda_v
        &= 1.
    \label{eq:peps_gaugeConsistency_connectivity_gap_global_scalar_product}
\end{align}
A spanning-tree construction then produces edge scalars whose oriented
incidence products realize the reciprocals:
\begin{align}
    \prod_{e\ni v}s_{v,e}
        &= \lambda_v^{-1}
        \qquad\text{for every }v\in V.
    \label{eq:peps_gaugeConsistency_connectivity_gap_connected_scalar_solution}
\end{align}
These scalars absorb the vertex proportionalities into one global gauge
family.

\section{What is established}

\begin{itemize}
    \item
    \emph{Per-vertex scalar.}
    Combining
    \leanid{TNLean.PEPS.post_absorption_edge_insertion_equality} with the four
    two-block injectivity facts
    \leanid{TNLean.PEPS.isTwoBlockInjective_vertexTwoBlock} and
    \leanid{TNLean.PEPS.isTwoBlockInjective_complementTwoBlock}, applied to
    $A$ and the reindexed $\widetilde B$, gives the required insertion
    hypotheses. The bridge
    \leanid{TNLean.PEPS.sameTwoBlockInsertions_of_edgeInsertedCoeff_eq} and
    \leanid{TNLean.PEPS.one_vertex_complement_comparison} then produce, at
    every non-isolated vertex $v$, a nonzero scalar $c_v$ satisfying
    \begin{align}
        A_v
            &= c_v\,\widetilde B_v.
        \label{eq:peps_gaugeConsistency_connectivity_gap_formal_vertex_scalar}
    \end{align}

    \item
    \emph{Gauge inversion.}
    Inverting the absorbed gauge gives
    \begin{align}
        B_v
            &= c_v^{-1}
               \operatorname{gaugeVertex}(A,Z^{-1})_v.
        \label{eq:peps_gaugeConsistency_connectivity_gap_inverse_absorbed_gauge}
    \end{align}
    This identity reduces gauge consistency to the edge-scalar incidence
    problem.

    \item
    \emph{Gauge assembly.}
    If $G$ is connected and
    \begin{align}
        \prod_{v\in V}c_v
            &= 1,
        \label{eq:peps_gaugeConsistency_connectivity_gap_formal_product_one}
    \end{align}
    then \leanid{TNLean.PEPS.exists_edgeScalars_of_connected} solves
    \begin{align}
        \prod_{e\ni v}s_{v,e}
            &= c_v^{-1}
            \qquad\text{for every }v\in V.
        \label{eq:peps_gaugeConsistency_connectivity_gap_formal_edge_scalar_solve}
    \end{align}
    Folding these scalars into the inverse absorbed gauge produces one global
    family satisfying
    \begin{align}
        B_v
            &= \operatorname{gaugeVertex}(A,X)_v
            \qquad\text{for every }v\in V.
        \label{eq:peps_gaugeConsistency_connectivity_gap_assembled_global_gauge}
    \end{align}
    This step is formalized by
    \leanid{TNLean.PEPS.perVertex_gauge_identity}. The cancellation
    \leanid{TNLean.PEPS.prod_perVertexScalar_eq_one} proves
    Eq.~\ref{eq:peps_gaugeConsistency_connectivity_gap_formal_product_one}.
    Together, these results reduce \leanid{TNLean.PEPS.gaugeConsistency} and
    \leanid{TNLean.PEPS.fundamentalTheorem_PEPS} to a single isolated lemma.

    \item
    \emph{State nonvanishing.}
    The product cancellation uses a nonzero state coefficient:
    \begin{align}
        \exists \sigma,\quad
        \langle A\rangle_\sigma
            &\neq 0.
        \label{eq:peps_gaugeConsistency_connectivity_gap_nonzero_state_coefficient}
    \end{align}
    This statement is formalized by
    \leanid{TNLean.PEPS.exists_stateCoeff_ne_zero}. For a vertex-injective PEPS
    with positive bond dimensions, it is the standing assumption in
    Ref.~\cite{Molnar2018NormalPEPS} that an injective PEPS is a genuine
    nonzero state.

    The \Lean{} derivation splits a state coefficient at one vertex against its
    complement using
    \leanid{TNLean.PEPS.stateCoeff_eq_vertexComplement}. The complement family
    is linearly independent by
    \leanid{TNLean.PEPS.vertexComplementTensorInjective_of_isVertexInjective}
    and therefore contains no zero member. Vertex injectivity at the selected
    vertex then forces a nonzero coefficient. Positive bond dimensions make
    the vertex-star configuration type nonempty. This argument requires no
    induced-subgraph or leaf-peeling infrastructure. It completes the
    reduction, so \leanid{TNLean.PEPS.fundamentalTheorem_PEPS} is fully proved
    and axiom-clean.
\end{itemize}

\section{Verdict}

The source claim is false as printed for disconnected graphs. The empty-graph
witness satisfies the stated injectivity and positive-bond hypotheses, has
equal PEPS states, and admits no global edge gauge relating the local tensors.
Adding connectivity gives the intended theorem, and the corrected formal
development is resolved.

This connectivity gap differs from the balanced-edge-scalar gap recorded in
\path{docs/paper-gaps/peps_gauge_edge_scalars.tex}. That note concerns the
uniqueness clause
\leanid{TNLean.PEPS.gauge_unique_mod_edge_scalars}, for which even a connected
graph has non-global edge-scalar freedom. The present note concerns existence:
disconnectedness prevents the assembly of one global gauge family.

\bibliographystyle{alpha}
\bibliography{references}

\end{document}
